% --- Tabla: Verificación vs. validación (Actividad 2) ---
\begin{table}[H]
  \centering
  \small
  \caption{Clasificación de actividades de aseguramiento de calidad entre verificación y validación.}
  \label{tab:pruebas_verificacion_validacion}
  \begin{tabularx}{\linewidth}{|L{4.6cm}|C{2.3cm}|C{2.0cm}|Y|}
    \hline
    \textbf{Acción} & \textbf{Verificación} & \textbf{Validación} & \textbf{Justificación} \\ \hline
    a) Revisar que el diagrama de clases respete el CUS «Registrar usuario». & \textbf{X} & &
      Se compara un producto de diseño (diagrama de clases) contra su especificación (el CUS). \\ \hline
    b) Mostrar el prototipo al cliente para confirmar que resuelve su necesidad. & & \textbf{X} &
      Se confirma con el cliente que el prototipo resuelve su necesidad real. \\ \hline
    c) Ejecutar las pruebas unitarias del módulo de inicio de sesión. & \textbf{X} & &
      Las pruebas unitarias comprueban que el código cumple su especificación. \\ \hline
    d) El usuario final prueba el sistema en su entorno y firma la conformidad. & & \textbf{X} &
      El usuario final confirma en su entorno que el sistema es el que necesitaba. \\ \hline
    e) Inspeccionar el código para confirmar que sigue el estándar de codificación. & \textbf{X} & &
      Se revisa el código contra un estándar definido, sin ejecutar el sistema. \\ \hline
    f) Comprobar que el sistema ensamblado cumple los requisitos funcionales y de rendimiento. & \textbf{X} & &
      Se comprueba el sistema ensamblado contra los requisitos especificados (RF y RNF). \\ \hline
  \end{tabularx}
  \par\vspace{0.2cm}
  {\raggedright \small \emph{Nota.} Elaboración propia (2026). Diferenciación metodológica entre verificación (conformidad contra especificación) y validación (adecuación a la necesidad del usuario final).\par}
\end{table}
